\chapter*{致读者}
\phantomsection
\addcontentsline{toc}{chapter}{致读者}

\begin{enumerate}
\def\labelenumi{\arabic{enumi}.}
\item
  本《数学原理》丛书从数学的开端讲起，并给出完整的证明。原则上，它并不要求读者具备任何专门的数学知识，而只要求对数学推理有一定的熟悉以及一定的抽象思维能力。不过，它特别是为那些至少已较好地掌握大学数学课程头一两年内容的读者而写的。
\item
  我们所选用的叙述方法是公理化的，并且通常从一般过渡到特殊。证明的需要使题材具有一种严格固定的顺序。因此，除非读者已经具备相当广博的数学知识，否则某些考虑的用处不会立即显现在他面前。
\item
  丛书分为若干卷，每卷又分为若干章。下面列出在法文版中已全部或部分出版的各卷。凡有英译本时，括号中给出相应的英文书名。在本卷中，引用一律指英文版（如果有的话），否则指法文版。
\end{enumerate}

Théorie des Ensembles (Theory of Sets) 记作 E (S)

Algèbre (Algebra) --- A (A)

Topologie Générale (General Topology) --- TG (GT)

Fonctions d'une Variable Réelle (Functions of a Real Variable) --- FVR (FRV)

Espaces Vectoriels Topologiques (Topological Vector Spaces) --- EVT (TVS)

Intégration (Integration) --- INT (INT)

(Commutative Algebra) --- AC (CA)

Variétés Différentielles et Analytiques --- VAR

(Lie Groups and Lie Algebras) --- LIE (LIE)

Théories Spectrales --- TS

在前六卷（按上述顺序）中，正文里的每一陈述都只假定已知那些已在同一章中讨论过的结果，或在按如下顺序排列的前面各章中讨论过的结果：S；A，第 I 至 III 章；GT，第 I 至 III 章；A，从第 IV 章起；GT，从第 IV 章起；FRV；TVS；INT。

从第七卷起，读者通常会在其中找到关于该卷与其他各卷之间逻辑关系的确切说明（前六卷总假定为已知）。

\begin{enumerate}
\def\labelenumi{\arabic{enumi}.}
\setcounter{enumi}{3}
\item
  不过，我们有时在正文中插入了一些例子，它们涉及读者可能已经知道、但在本丛书中尚未讨论的事实。这类例子置于两个星号之间：*\ldots*。大多数读者无疑会发现这些例子有助于理解正文。在另一些情形，*\ldots* 之间的段落所涉及的结果在本丛书的别处讨论。我们希望读者能够自行验证其中不存在任何恶性循环。
\item
  每一章的逻辑框架由该章的定义、公理和定理组成。这些是以后使用时主要应当记住的部分。较次要的结果以及那些容易由定理推出的结果，标以“命题”、“引理”、“推论”、“注记”等名称。初读时可以略去的那些部分以小号字排印。对特别重要的定理所作的评注，偶尔以“scholium（评注）”之名出现。
\end{enumerate}

为了避免令人生厌的重复，有时方便的做法是引入仅在某一章或某一章的某一节内有效的记号或缩写（例如，在只讨论交换环的一章中，“环”一词总是指“交换环”）。这类约定总会被明确指出，一般是在出现该约定的一章的开头。

\pandocbounded{\includegraphics[keepaspectratio]{images/part001_c576909a.jpg}}

\begin{enumerate}
\def\labelenumi{\arabic{enumi}.}
\setcounter{enumi}{5}
\item
  正文中有一些段落旨在预先提醒读者谨防严重错误。这些段落用（“危险弯道”）记号在页边标出。
\item
  习题的目的，一方面是使读者能确信自己已经消化了正文，另一方面是使读者注意到那些在正文中没有位置却仍有意义的结果。最难的习题带有记号 ¶。
\item
  一般说来，我们遵循通行的术语，除非看来有充分理由偏离它。
\item
  我们特别致力于始终使用严格正确的语言而不牺牲简洁性。对于各种语言滥用，我们在正文中尽可能加以提醒；没有这类滥用，任何数学文本都有流于迂腐、甚至无法阅读的危险。
\item
  由于正文原则上是对一个理论的教条式陈述，它一般不包含对文献的引用。书目式引用汇集于历史注记之中。每篇历史注记后面所附的参考文献一般只包括那些对所述理论的演进曾起过最重要作用的书籍和原始论文。它绝不冒称在任何意义上是完备的。
\end{enumerate}

至于习题，我们一般不认为值得注明其出处，因为它们取自许多不同的来源（原始论文、教科书、习题集）。

\begin{enumerate}
\def\labelenumi{\arabic{enumi}.}
\setcounter{enumi}{10}
\tightlist
\item
  对本丛书某一部分的引用按如下方式给出：
\end{enumerate}

\begin{enumerate}
\def\labelenumi{\alph{enumi})}
\item
  若引用的是同一节（§）中陈述的定理、公理或定义，则按其编号引用。
\item
  若它们出现在同一章的另一节中，则在引用时同时注明该节。
\item
  若它们出现在同一卷的另一章中，则注明章与节。
\item
  若它们出现在另一卷中，则先注明该卷，用其书名的缩写表示。
\end{enumerate}

结果汇编以字母 R 引用；于是 S, R 表示“集合论的结果汇编”。

